%!TEX root = ../../document/document.tex

% For every chapter create a file in the /content/chapters folder with the name
% chapter-XX.tex where XX is the chapter number  (e.g. 01, 02, 03, ..., 99).


\chapter{Stand der Technik}\label{c:2}

\section{Python Flask}\label{c:2_s:flask}

Flask ist ein leichtgewichtiges Framework für Python, das eine einfache und flexible Möglichkeit bietet, Webanwendungen zu erstellen. Es ist bekannt für seine Einfachheit und Modularität, was es Entwicklern ermöglicht, schnell Prototypen zu erstellen. Andererseits ist es ebenfalls möglich, im Laufe der Entwicklung die Anwendungen zu skalieren. Flask folgt dem \ac{WSGI}-Standard und ist in der Lage, mit verschiedenen Webservern zu kommunizieren.
